\documentclass[12pt,a4paper]{article}
\usepackage[margin=0.8in]{geometry}
\usepackage{amsmath,amssymb}
\usepackage{enumitem}
\usepackage{titlesec}
\usepackage{fancyhdr}

\pagestyle{fancy}
\fancyhf{}
\rhead{Quadratic Equations -- Class XI}
\lhead{All Questions with Final Answers}
\cfoot{\thepage}

\titleformat{\section}{\Large\bfseries}{\thesection}{1em}{}
\titleformat{\subsection}{\large\bfseries}{\thesubsection}{1em}{}

\title{\textbf{Quadratic Equations -- Class XI}\\[0.4cm]
\large Complete Extraction: All Questions + Final Answers Only}
\author{}
\date{}

\begin{document}

\maketitle
\tableofcontents
\newpage

%============================================================
\section{5.1 Solving Quadratic Equations by Formula}
%============================================================

\subsection{Illustrative Examples}

\begin{enumerate}[label=Example \arabic*.]
    \item Solve the following quadratic equations:
    \begin{enumerate}[label=(\roman*)]
        \item \( 3x^2 - 4x - 4 = 0 \)
        \item \( 4x^2 - 2x + \dfrac{1}{4} = 0 \)
        \item \( 3x^2 - 4x + \dfrac{20}{3} = 0 \)
    \end{enumerate}

    \item Solve the following equations:
    \begin{enumerate}[label=(\roman*)]
        \item \( x^2 + 3x + 9 = 0 \)
        \item \( 27x^2 - 10x + 1 = 0 \)
    \end{enumerate}

    \item Solve the following equations:
    \begin{enumerate}[label=(\roman*)]
        \item \( \sqrt{3}x^2 - \sqrt{2}x + 3\sqrt{3} = 0 \)
        \item \( x^2 + \dfrac{x}{\sqrt{2}} + 1 = 0 \)
    \end{enumerate}

    \item Find the roots of the equation \( x^2 + x - (a + 2)(a + 1) = 0 \).

    \item Solve the equation \( 2x^2 - 5x + 2 = 0 \) by the method of completing the squares.

    \item Find two numbers such that their sum is 6 and the product is 14.

    \item Three consecutive natural numbers are such that the square of the middle number exceeds the difference of the squares of the other two by 60. Assuming the middle number to be \( x \), form a quadratic equation and find the three numbers.

    \item A two digit number is four times the sum and twice the product of its digits. Find the number.

    \item A factory kept increasing its output by the same percentage every year. If its output doubled in the last two years, find the percentage increase every year.
\end{enumerate}

\subsection{Answers -- Illustrative Examples (5.1)}

\begin{enumerate}[label=Example \arabic*.]
    \item 
    \begin{enumerate}[label=(\roman*)]
        \item \( 2, -\dfrac{2}{3} \)
        \item \( \dfrac{1}{4}, \dfrac{1}{4} \)
        \item \( \dfrac{2 \pm 4i}{3} \)
    \end{enumerate}

    \item 
    \begin{enumerate}[label=(\roman*)]
        \item \( \dfrac{-3 \pm 3\sqrt{3}i}{2} \)
        \item \( \dfrac{5 \pm \sqrt{2}i}{27} \)
    \end{enumerate}

    \item 
    \begin{enumerate}[label=(\roman*)]
        \item \( \dfrac{\sqrt{2} \pm \sqrt{34}i}{2\sqrt{3}} \)
        \item \( \dfrac{-1 \pm \sqrt{7}i}{2\sqrt{2}} \)
    \end{enumerate}

    \item \( a + 1,\ -(a + 2) \)

    \item \( 2,\ \dfrac{1}{2} \)

    \item \( 3 + \sqrt{5}i,\ 3 - \sqrt{5}i \)

    \item \( 9,\ 10,\ 11 \)

    \item \( 36 \)

    \item \( 100(\sqrt{2} - 1)\% \)
\end{enumerate}

%------------------------------------------------------------
\subsection{Exercise 5.1}

\textbf{Very short answer / objective questions (1 to 5)}

Solve the following quadratic equations:

\begin{enumerate}
    \item (i) \( x^2 + 2 = 0 \) \quad (ii) \( 4x^2 + 5 = 0 \)
    \item (i) \( x^2 - x - 12 = 0 \) \quad (ii) \( 25x^2 + 30x + 7 = 0 \)
    \item (i) \( \sqrt{3}x^2 - 4x + \sqrt{3} = 0 \) \quad (ii) \( 15x^2 + 15 = 34x \)
    \item (i) \( x^2 + x + 1 = 0 \) \quad (ii) \( x^2 - x + 2 = 0 \)
    \item (i) \( x^2 + 3x + 5 = 0 \) \quad (ii) \( 2x^2 + x + 1 = 0 \)
\end{enumerate}

\textbf{Short answer questions (6 to 8)}

\begin{enumerate}[start=6]
    \item (i) \( \sqrt{2}x^2 + x + \sqrt{2} = 0 \) \quad (ii) \( \sqrt{5}x^2 + x + \sqrt{5} = 0 \)
    \item (i) \( 21x^2 - 28x + 10 = 0 \) \quad (ii) \( x^2 - 2x + \dfrac{3}{2} = 0 \)
    \item Find the values of \( x \) if \( p + 7 = 0 \), \( q - 12 = 0 \) and \( x^2 + px + q = 0 \).
\end{enumerate}

\textbf{Long answer questions (9 to 17)}

\begin{enumerate}[start=9]
    \item Divide 39 into 2 parts whose product is 338.
    \item The sum of two numbers is 9 and the sum of their squares is 41. Taking one number as \( x \), form an equation in \( x \) and solve it.
    \item The sum of the squares of two consecutive positive even integers is 340. Find the integers.
    \item A two digit number is four times the sum and three times the product of its digits. Find the number.
    \item The hypotenuse of a right angled triangle is 17 cm long and the difference between the other two sides is 7 cm. Find the lengths of the two unknown sides.
    \item Originally a rectangle was twice as long as it is wide. When 4 m were added to its length and 3 m subtracted from its width, the resulting rectangle had an area of 600 m\(^2\). Find the dimensions of the new rectangle.
    \item A rectangular garden 10 m by 16 m is to be surrounded by a concrete walk of uniform width. Given that area of the walk is 120 square metres, find its width.
    \item Two years ago, a man's age was three times the square of his son's age. In three years time, his age will be four times his son's age. Find their present ages.
    \item A piece of cloth costs Rs175. If the piece were 4 m longer and each metre costs Rs5 less, the cost would have remained unchanged. How long is the piece?
\end{enumerate}

\subsection{Answers -- Exercise 5.1 (Final Answers Only)}

\begin{enumerate}
    \item (i) \( \pm \sqrt{2}i \) \quad (ii) \( \pm \dfrac{\sqrt{5}}{2}i \)
    \item (i) \( 4, -3 \) \quad (ii) \( -\dfrac{3 + \sqrt{2}}{5} \)
    \item (i) \( \sqrt{3},\ \dfrac{1}{\sqrt{3}} \) \quad (ii) \( \dfrac{3}{5},\ \dfrac{5}{3} \)
    \item (i) \( -\dfrac{1 + \sqrt{5}i}{2} \) \quad (ii) \( \dfrac{1 \pm \sqrt{7}i}{2} \)
    \item (i) \( -\dfrac{3 \pm \sqrt{11}i}{2} \) \quad (ii) \( -\dfrac{1 + \sqrt{7}i}{4} \)
    \item (i) \( -\dfrac{1 + \sqrt{7}i}{2\sqrt{2}} \) \quad (ii) \( -\dfrac{1 \pm \sqrt{19}i}{2\sqrt{5}} \)
    \item (i) \( \dfrac{14 \pm \sqrt{14}i}{21} \) \quad (ii) \( \dfrac{2 \pm \sqrt{2}i}{2} \)
    \item \( 3, 4 \)
    \item \( 26, 13 \)
    \item \( 4, 5 \)
    \item \( 12, 14 \)
    \item \( 24 \)
    \item \( 8 \) cm, \( 15 \) cm
    \item Length = 40 m, Width = 15 m
    \item 2 metres
    \item 29 years, 5 years
    \item 10 metres
\end{enumerate}

%============================================================
\section{5.2 Equations Reducible to Quadratic Form}
%============================================================

\subsection{Illustrative Examples}

\begin{enumerate}[label=Example \arabic*.]
    \item Find the roots of the equation \( \dfrac{1}{x+1} + \dfrac{2}{x+2} = \dfrac{4}{x+4} \).

    \item Solve: \( \dfrac{x-1}{x-2} - \dfrac{x-2}{x-3} = \dfrac{x-5}{x-6} - \dfrac{x-6}{x-7} \).

    \item Solve: \( \dfrac{a}{ax+1} + \dfrac{b}{bx+1} = a + b \), \( ab \neq 0 \), \( a + b \neq 0 \).

    \item Solve: \( 2^{2x+8} + 1 = 32 \cdot 2^x \).

    \item Solve: \( (x^2 - 5x)^2 - 30(x^2 - 5x) - 216 = 0 \).

    \item Solve: \( (x+1)(x+2)(x+3)(x+4) = 120 \).

    \item Solve: \( x^{2/3} + x^{1/3} = 2 \).

    \item Solve the equation \( 8 \sqrt{\dfrac{x}{x+3}} - \sqrt{\dfrac{x+3}{x}} = 2 \).

    \item Solve the equation \( \sqrt{a+x} + \sqrt{a-x} = \dfrac{a}{x} \), \( x \neq 0 \).

    \item Solve \( \sqrt{x} = x - 2 \).

    \item Solve the equation \( \sqrt{x^2 + 3x + 32} + \sqrt{x^2 + 3x + 5} = 9 \).

    \item Solve the equation \( \sqrt{x+5} + \sqrt{x+21} = \sqrt{6x+40} \).

    \item Solve the equation \( \left(x + \dfrac{1}{x}\right)^2 = 4 + \dfrac{3}{2}\left(x - \dfrac{1}{x}\right) \).

    \item Solve \( 4x^4 - 16x^3 + 7x^2 + 16x + 4 = 0 \).

    \item Solve: \( (x + 2)(x + 3)(x + 8)(x + 12) = 4x^2 \).
\end{enumerate}

\subsection{Answers -- Illustrative Examples (5.2)}

\begin{enumerate}[label=Example \arabic*.]
    \item \( 2 + 2\sqrt{3},\ 2 - 2\sqrt{3} \)
    \item \( 4\frac{1}{2} \)
    \item \( 0,\ -\dfrac{a^2 + b^2}{ab(a + b)} \)
    \item \( x = -4 \)
    \item \( 9,\ -4,\ 3,\ 2 \)
    \item \( 1,\ -6,\ \dfrac{-5 \pm i\sqrt{39}}{2} \)
    \item \( 1,\ -8 \)
    \item \( x = 1 \)
    \item \( a,\ -a \)
    \item \( 4 \)
    \item \( -4,\ 1 \)
    \item \( 4 \)
    \item \( -\dfrac{1}{2},\ -1,\ 1,\ 2 \)
    \item \( 2,\ -\dfrac{1}{2},\ \dfrac{5 \pm \sqrt{41}}{4} \)
    \item \( -4,\ -6,\ \dfrac{-15 \pm \sqrt{129}}{2} \)
\end{enumerate}

%------------------------------------------------------------
\subsection{Exercise 5.2}

Solve the following equations:

\begin{enumerate}
    \item (i) \( \dfrac{x-2}{x+2} + \dfrac{x+2}{x-2} = 4 \)
    \item (ii) \( \dfrac{x-1}{x-2} + \dfrac{x-3}{x-4} = 3\frac{1}{3} \)
    \item \( \dfrac{1}{x+1} - \dfrac{2}{x+2} = \dfrac{3}{x+3} - \dfrac{4}{x+4} \)
    \item (i) \( 3^x + 3^{-x} = 2 \) \quad (ii) \( 2^{x+1} + 4^x = 8 \)
    \item \( 5^x + 1 + 5^{2-x} = 5^3 + 1 \)
    \item \( x^{2/3} - x^{1/3} - 2 = 0 \)
    \item (i) \( (x^2 + 3x + 2)^2 - 8(x^2 + 3x) - 4 = 0 \)
    \item (ii) \( (x^2 - 5x + 7)^2 - (x-2)(x-3) = 1 \)
    \item \( x(x+2)(x+3)(x+5) = 72 \)
    \item \( x(x+2)(x^2 - 1) = -1 \)
    \item \( (x+1)(2x+3)(2x+5)(x+3) = 945 \)
    \item \( (x^2 + x - 6)(x^2 - 3x - 4) = 24 \)
    \item (i) \( \sqrt{\dfrac{x}{1-x}} + \sqrt{\dfrac{1-x}{x}} = \dfrac{13}{6} \)
    \item (ii) \( \sqrt{\dfrac{4x-1}{4x+1}} - \sqrt{\dfrac{4x+1}{4x-1}} = \dfrac{8}{3} \), \( x \neq \pm \dfrac{1}{4} \)
    \item (i) \( \sqrt{3x+4} = x \) \quad (ii) \( \sqrt{3x+1} - \sqrt{x-1} = 2 \)
    \item (i) \( \sqrt{x+1} + \sqrt{x-1} = 1 \) \quad (ii) \( \sqrt{x+3} + \sqrt{x-3} = 6 \)
    \item (i) \( \sqrt{3x^2 - 4x + 34} - \sqrt{3x^2 - 4x - 11} = 5 \)
    \item (ii) \( \sqrt{3x^2 - 7x - 30} - \sqrt{2x^2 - 7x - 5} = x - 5 \)
    \item \( 12 + 9\sqrt{(x-1)(3x+2)} = 3x^2 - x \)
    \item \( \dfrac{\sqrt{x^2+1} - \sqrt{x^2-1}}{\sqrt{x^2+1} + \sqrt{x^2-1}} = \dfrac{1}{3} \)
    \item (i) \( \left(x + \dfrac{1}{x}\right)^2 - 2\left(x - \dfrac{1}{x} + 4\right) = 11 \)
    \item (ii) \( 12x^4 - 56x^3 + 89x^2 - 56x + 12 = 0 \)
    \item Solve \( (x^2 + 2)^2 + 8x^2 = 6x(x^2 + 2) \)
\end{enumerate}

\subsection{Answers -- Exercise 5.2 (Final Answers Only)}

\begin{enumerate}
    \item (i) \( 2\sqrt{3},\ -2\sqrt{3} \) \quad (ii) \( 5,\ \dfrac{5}{2} \)
    \item \( 0,\ -\dfrac{5}{2} \)
    \item (i) \( 0 \) \quad (ii) \( 1 \)
    \item \( -1,\ 2 \)
    \item \( -1,\ 8 \)
    \item (i) \( -4,\ -3,\ 0,\ 1 \) \quad (ii) \( 2,\ 3,\ \dfrac{5 \pm \sqrt{3}i}{2} \)
    \item \( 1,\ -6,\ \dfrac{-5 \pm i\sqrt{23}}{2} \)
    \item \( -\dfrac{1 + \sqrt{5}}{2},\ -\dfrac{1 - \sqrt{5}}{2} \)
    \item \( 2,\ -6,\ \dfrac{-4 + i\sqrt{59}}{2} \)
    \item \( 0,\ 1,\ \dfrac{1 + \sqrt{57}}{2} \)
    \item (i) \( \dfrac{9}{13},\ \dfrac{4}{13} \) \quad (ii) \( -\dfrac{5}{16} \)
    \item (i) \( 4 \) \quad (ii) \( 1,\ 5 \)
    \item (i) No solution \quad (ii) \( \dfrac{37}{4} \)
    \item (i) \( 3,\ -\dfrac{5}{3} \) \quad (ii) \( 5,\ 6 \)
    \item \( 6,\ -\dfrac{17}{3} \)
    \item \( \pm \sqrt{\dfrac{5}{3}} \)
    \item (i) \( \dfrac{5 \pm \sqrt{29}}{2},\ -\dfrac{3 \pm \sqrt{13}}{2} \) \quad (ii) \( \dfrac{1}{2},\ \dfrac{2}{3},\ \dfrac{3}{2},\ 2 \)
    \item \( 1 \pm i,\ 2 \pm \sqrt{2} \)
\end{enumerate}

%============================================================
\section{5.3 Nature of Roots of a Quadratic Equation}
%============================================================

\subsection{Illustrative Examples}

\begin{enumerate}[label=Example \arabic*.]
    \item Discuss the nature of the roots of the following equations:
    \begin{enumerate}[label=(\roman*)]
        \item \( 4x^2 - 12x + 9 = 0 \)
        \item \( 3x^2 - 10x + 3 = 0 \)
        \item \( 9x^2 - 2 = 0 \)
        \item \( x^2 + x + 1 = 0 \)
    \end{enumerate}

    \item Discuss the nature of the roots of the following equations:
    \begin{enumerate}[label=(\roman*)]
        \item \( 4x^2 - 4\sqrt{3}x + 3 = 0 \)
        \item \( 2\sqrt{3}x^2 - 5x + \sqrt{3} = 0 \)
    \end{enumerate}

    \item If \( a, b, h \in \mathbb{R} \), show that the equation \( (x - a)(x - b) = h^2 \) has real roots.

    \item If the roots of \( ax^2 + x + b = 0 \) are real and distinct, show that the roots of the equation \( \dfrac{x^2 + 1}{x} = 4\sqrt{ab} \) are imaginary.

    \item Given that \( a, b, c \in \mathbb{Q} \), show that the roots of the equation \( (b - c)x^2 + (c - a)x + (a - b) = 0 \) are rational. Find the condition for the roots to be equal.

    \item Discuss the nature of the roots of the equation \( (m + 6)x^2 + (m + 6)x + 2 = 0 \).

    \item For what value of \( k \) does the expression \( x^2 - 2x(1 + 3k) + 7(3 + 2k) \) become a perfect square?

    \item Determine a positive real value of \( k \) such that both the equations \( x^2 + kx + 64 = 0 \) and \( x^2 - 8x + k = 0 \) may have real roots.

    \item Show that if \( p, q, r, s \) are real numbers and \( pr = 2(q + s) \), then at least one of the equations \( x^2 + px + q = 0 \) and \( x^2 + rx + s = 0 \) has real roots.
\end{enumerate}

\subsection{Answers -- Illustrative Examples (5.3)}

\begin{enumerate}[label=Example \arabic*.]
    \item 
    \begin{enumerate}[label=(\roman*)]
        \item Rational and equal
        \item Rational and unequal
        \item Irrational and unequal
        \item Pair of complex conjugates
    \end{enumerate}

    \item 
    \begin{enumerate}[label=(\roman*)]
        \item Real and equal
        \item Real and unequal
    \end{enumerate}

    \item Proven (discriminant \( \geq 0 \))

    \item Proven (discriminant of second equation is negative)

    \item Roots are rational. Equal when \( b = \dfrac{a + c}{2} \)

    \item 
    \begin{itemize}
        \item Real \& equal when \( m = 2 \) or \( m = -6 \)
        \item Real \& unequal when \( m < -6 \) or \( m > 2 \)
        \item Complex when \( -6 < m < 2 \)
    \end{itemize}

    \item \( k = 2 \) or \( k = -\dfrac{10}{9} \)

    \item \( k = 16 \)

    \item Proven
\end{enumerate}

%------------------------------------------------------------
\subsection{Exercise 5.3}

\begin{enumerate}
    \item Find the nature of roots of the following equations without solving them:
    \begin{enumerate}[label=(\roman*)]
        \item \( x^2 + 9 = 0 \)
        \item \( 4x^2 - 24x + 35 = 0 \)
        \item \( x^2 - 2\sqrt{2}x + 1 = 0 \)
        \item \( 2x^2 - 2\sqrt{5}x + 3 = 0 \)
    \end{enumerate}

    \item (i) Prove that the roots of the equation \( x^2 - 2ax + a^2 - b^2 - c^2 = 0 \) are always real.

    \item (ii) Show that roots of the equation \( (x - a)(x - b) = abx^2 \) are always real. When are they equal?

    \item Show that the roots of the equation \( (x - a)(x - b) + (x - b)(x - c) + (x - c)(x - a) = 0 \) are always real. Find the condition that the roots may be equal.

    \item Discuss the nature of roots of the following equations:
    \begin{enumerate}[label=(\roman*)]
        \item \( \sqrt{3}x^2 - 2x - \sqrt{3} = 0 \)
        \item \( x^2 - (p + 1)x + p = 0 \)
        \item \( (x - a)(x - b) = ab \)
    \end{enumerate}

    \item Find \( m \) so that roots of the equation \( (4 + m)x^2 + (m + 1)x + 1 = 0 \) may be equal.

    \item Show that the roots of the equation \( x^2 + 2(3a + 5)x + 2(9a^2 + 25) = 0 \) are complex unless \( a = \dfrac{5}{3} \).

    \item If \( a, b, c, d \in \mathbb{R} \), show that the roots of the equation \( (a^2 + c^2)x^2 + 2(ab + cd)x + (b^2 + d^2) = 0 \) cannot be real unless they are equal.

    \item If \( a, b, c \in \mathbb{Q} \) and \( a + b + c = 0 \), show that roots of the equation \( ax^2 + bx + c = 0 \) are rational.

    \item Prove that the equation \( x^2 + px - 1 = 0 \) has real and distinct roots for all real values of \( p \).

    \item Show that the roots of the equation \( x^2 - 2\left( \dfrac{m+1}{m} \right)x + 3 = 0 \) are real for all (non-zero) real values of \( m \).

    \item If the roots of the equation \( ax^2 + 2cx + b = 0 \) are real and distinct, then show that the roots of the equation \( x^2 - 2(a + b)x + a^2 + b^2 + 2c^2 = 0 \) are non-real complex numbers.

    \item If the roots of the equation \( a(b - c)x^2 + b(c - a)x + c(a - b) = 0 \) are equal, show that \( \dfrac{2}{b} = \dfrac{1}{a} + \dfrac{1}{c} \).

    \item For what value of \( k \), \( (4 - k)x^2 + 2(k + 2)x + (8k + 1) \) is a perfect square.

    \item If the roots of the equation \( x^2 - 8x + m(m - 6) = 0 \) are real and distinct, then find all possible values of \( m \).
\end{enumerate}

\subsection{Answers -- Exercise 5.3 (Final Answers Only)}

\begin{enumerate}
    \item (i) Pair of complex conjugates
    \item (ii) Rational and unequal
    \item (iii) Real and unequal
    \item (iv) Pair of complex conjugates
    \item (ii) Roots are real and equal when \( a = b = 0 \)
    \item Roots are equal when \( a = b = c \). Then roots are \( a, a \)
    \item (i) Real and distinct
    \item (ii) Rational and distinct when \( p \neq 1 \); equal when \( p = 1 \)
    \item (iii) Real and distinct when \( a + b \neq 0 \); both zero when \( a + b = 0 \)
    \item \( m = 5 \) or \( m = -3 \)
    \item Proven
    \item Proven
    \item Proven
    \item Proven
    \item Proven
    \item Proven
    \item \( k = 0 \) or \( k = 3 \)
    \item \( -2 < m < 8 \)
\end{enumerate}

%============================================================
\section{5.4 Relations between Roots and Coefficients}
%============================================================

\subsection{Illustrative Examples}

\begin{enumerate}[label=Example \arabic*.]
    \item If \( \alpha, \beta \) are roots of the equation \( x^2 - 4x + 2 = 0 \), find the values of:
    \begin{enumerate}[label=(\roman*)]
        \item \( \alpha^2 + \beta^2 \)
        \item \( \alpha^2 - \beta^2 \)
        \item \( \alpha^3 + \beta^3 \)
        \item \( \dfrac{1}{\alpha} + \dfrac{1}{\beta} \)
    \end{enumerate}

    \item If \( \alpha, \beta \) are the roots of the equation \( ax^2 + bx + c = 0 \), then show that \( ax^2 + bx + c = a(x - \alpha)(x - \beta) \).

    \item For the quadratic equation \( (k-1)x^2 = kx - 1 \), \( k \neq 1 \), find \( k \) so that:
    \begin{enumerate}[label=(\roman*)]
        \item one root is \( -3 \)
        \item the sum of the roots is 2
        \item the product of the roots is \( -3 \)
        \item the roots are equal
        \item the roots are numerically equal but opposite in sign
    \end{enumerate}

    \item If one of the roots of \( 3x^2 - 7px + 10p = 0 \) be 5, find \( p \) and also the other root.

    \item The sum of the roots of the equation \( \dfrac{1}{x+a} + \dfrac{1}{x+b} = \dfrac{1}{c} \) is zero. Prove that the product of the roots is \( -\dfrac{1}{2}(a^2 + b^2) \).

    \item If \( \alpha, \beta \) be the roots of the equation \( ax^2 + bx + c = 0 \), evaluate:
    \begin{enumerate}[label=(\roman*)]
        \item \( \alpha^3\beta + \alpha\beta^3 \)
        \item \( \dfrac{\alpha^2}{\beta} + \dfrac{\beta^2}{\alpha} \)
        \item \( (a\alpha + b)^2 + (a\beta + b)^2 \)
    \end{enumerate}
\end{enumerate}

\subsection{Answers -- Illustrative Examples (5.4)}

\begin{enumerate}[label=Example \arabic*.]
    \item 
    \begin{enumerate}[label=(\roman*)]
        \item \( 12 \)
        \item \( \pm 8\sqrt{2} \)
        \item \( 40 \)
        \item \( 2 \)
    \end{enumerate}

    \item Proven

    \item 
    \begin{enumerate}[label=(\roman*)]
        \item \( k = \dfrac{2}{3} \)
        \item \( k = 2 \)
        \item \( k = \dfrac{2}{3} \)
        \item \( k = 2 \)
        \item \( k = 0 \)
    \end{enumerate}

    \item \( p = 3 \), other root = 2

    \item Proven

    \item 
    \begin{enumerate}[label=(\roman*)]
        \item \( \dfrac{c(b^2 - 2ac)}{a^3} \)
        \item \( \dfrac{b(b^2 - 2ac)}{a^2c} \)
        \item \( \dfrac{b^2(b^2 - 2ac)}{a^2} \)
    \end{enumerate}
\end{enumerate}

\end{document}
